\chapter{Development} \label{cap:development}

The rvsc0 and rvsc1 targets restrict the instruction set available to the compiler
(Chapter~\ref{cap:requirements} specifies the complete native instruction set for each target). Every
C construct the compiler may emit must be realized using only those native instructions; for every
excluded instruction, GCC synthesizes an equivalent sequence at compile time. This chapter presents
the tools and infrastructure used, the mathematical derivation and proof of each synthesis, the GCC
implementation that makes the process transparent to the programmer, and the known limitations of each
target.

\section{Technologies Used}

\subsection{GCC 17.0.0}

The custom targets are implemented as a backend extension to GCC 17.0.0. GCC is chosen because it is
the dominant open-source compiler for embedded and systems software and because its machine description
framework provides a declarative mechanism for defining instruction synthesis rules that is already
present in the upstream RISC-V backend. The following files in the GCC source tree were modified or
added:

\begin{table}[htb]
\centering
\caption{GCC source files modified or added for this project}
\label{tab:gcc-files}
\begin{tabular}{ll}
\toprule
\textbf{File} & \textbf{Role} \\
\midrule
\texttt{gcc/config/riscv/riscv.md}  & Machine description: synthesis patterns and native instruction guards \\
\texttt{gcc/config/riscv/riscv.opt} & Option declarations: per-synthesis boolean flags \\
\texttt{gcc/config/riscv/riscv.cc}  & Target hooks: addi/shift constant synthesis for LUI \\
\texttt{gcc/config/riscv/rvscN.h}   & Per-target headers: \texttt{CC1\_SPEC} injecting \texttt{-mno-*} flags automatically \\
\texttt{gcc/config/config.gcc}      & Triple mapping: \texttt{rvscN-*-elf*} to \texttt{cpu\_type=riscv} \\
\texttt{gcc/config/config.sub}      & Triple normalisation: recognises \texttt{rvscN} as a valid CPU name \\
\bottomrule
\end{tabular}
\end{table}

\subsection{GNU Binutils (riscv32/64-none-elf)}

The upstream GNU binutils cross toolchain is used without modification as the assembler and linker.
The binutils triple differs from the compiler triple: \texttt{riscv32-none-elf} for rvsc0--rvsc3 and
\texttt{riscv64-none-elf} for rvsc4--rvsc7.

\subsection{Spike 1.1.1-dev}

Spike is the official RISC-V ISA reference simulator \cite{spike}. It implements the full
RV64IMAFD instruction set, making it suitable as the behavioral oracle: a binary compiled by
\texttt{rvsc1-unknown-elf-gcc} contains only valid RV32I instructions (the synthesis sequences are
themselves valid RV32I), so Spike can execute it and report the correct result. Spike is invoked with
\texttt{--isa=rv32i} for rvsc0--rvsc3 tests, and with \texttt{--log-commits} to count retired
instructions for performance measurements.

\section{Toolchain Installation}

\subsection{Repository Setup}

The GCC backend modifications live in a fork of the upstream GCC repository:

\begin{lstlisting}[language=bash]
git clone https://github.com/guissalustiano/gcc-hannersy-paterson gcc
\end{lstlisting}

The GNU Binutils source is cloned separately and is required for all targets, since GCC must invoke
the target assembler during compilation:

\begin{lstlisting}[language=bash]
git clone https://sourceware.org/git/binutils-gdb.git binutils-gdb
\end{lstlisting}

Both repositories should sit side by side in the same parent directory. The build directories and
installed toolchains are placed inside a per-target subdirectory:

\begin{lstlisting}
project/
+-- gcc/              <- GCC fork (source, never built in-tree)
+-- binutils-gdb/     <- binutils source
+-- targets/
    +-- rvsc2/
        +-- build-binutils/   <- binutils out-of-tree build
        +-- build/            <- GCC out-of-tree build
        +-- install/          <- installed toolchain (bin/, lib/, ...)
\end{lstlisting}

\subsection{Building Binutils}

Binutils must be built and installed before GCC so that GCC's configure step can detect the target
assembler and linker.

\begin{lstlisting}[language=bash]
mkdir -p targets/rvsc2/build-binutils
cd targets/rvsc2/build-binutils

../../binutils-gdb/configure \
    --target=rvsc2-unknown-elf \
    --prefix=$(pwd)/../install \
    --disable-nls \
    --disable-gdb

make -j$(nproc)
make install
\end{lstlisting}

After this step the install prefix contains the target assembler, linker, and binutils utilities
(\texttt{rvsc2-unknown-elf-as}, \texttt{rvsc2-unknown-elf-ld}, etc.).

\subsection{Building GCC}

With the assembler and linker installed, GCC can be configured and built. The three
\texttt{--disable-*} flags prevent GCC from attempting to rebuild its own copies of the binutils
components from the source tree (the GCC repository bundles a copy of binutils), which avoids
duplicate build failures and unnecessary compilation time.

\begin{lstlisting}[language=bash]
mkdir -p targets/rvsc2/build
cd targets/rvsc2/build

PATH="$(pwd)/../install/bin:$PATH" \
../../gcc/configure \
    --target=rvsc2-unknown-elf \
    --prefix=$(pwd)/../install \
    --enable-languages=c \
    --with-newlib \
    --disable-binutils \
    --disable-ld \
    --disable-gas

make all-gcc -j$(nproc)
make install-gcc
\end{lstlisting}

Prepending the install prefix to \texttt{PATH} before configure allows the configure script to detect
the pre-installed \texttt{rvsc2-unknown-elf-as} and \texttt{rvsc2-unknown-elf-ld} binaries, preventing
it from scheduling them for in-tree compilation.

\subsection{Building Newlib and libgcc}

With GCC installed, the C runtime library (newlib) and the compiler support library (libgcc) can be
built. Both are required to link executable programs.

\begin{lstlisting}[language=bash]
mkdir -p targets/rvsc2/build-newlib
cd targets/rvsc2/build-newlib

PATH="$(pwd)/../install/bin:$PATH" \
../../newlib-src/configure \
    --target=rvsc2-unknown-elf \
    --prefix=$(pwd)/../install \
    --disable-multilib \
    --disable-newlib-supplied-syscalls

PATH="$(pwd)/../install/bin:$PATH" make -j$(nproc)
PATH="$(pwd)/../install/bin:$PATH" make install
\end{lstlisting}

With newlib installed, libgcc can be built against it:

\begin{lstlisting}[language=bash]
cd targets/rvsc2/build

PATH="$(pwd)/../install/bin:$PATH" make all-target-libgcc -j$(nproc)
PATH="$(pwd)/../install/bin:$PATH" make install-target-libgcc
\end{lstlisting}

After all four steps the toolchain is fully self-contained under \texttt{targets/rvsc2/install/bin/}.
Compiling a C program to a RISC-V object file requires only:

\begin{lstlisting}[language=bash]
PATH="targets/rvsc2/install/bin:$PATH" \
rvsc2-unknown-elf-gcc -O2 -ffreestanding -c program.c -o program.o
\end{lstlisting}

\section{Synthesis Derivations} \label{sec:sc1-synthesis}

\textbf{Assembly notation.} In the listings below, \texttt{rd}, \texttt{rs1}, and \texttt{rs2} denote
the canonical destination and source registers. Registers \texttt{t0}--\texttt{t5} are temporaries
chosen for readability. In the actual GCC machine description, all temporaries are allocated as
pseudo-registers via \texttt{gen\_reg\_rtx(SImode)}. The register allocator maps them to physical
registers, resolving aliasing conflicts automatically. The bracket notation \texttt{[op ...]} marks
an instruction that is itself synthesized, its expansion is defined in the subsection that covers that
operation.

\subsection{Arithmetic}

Arithmetic synthesis operations assign values only to registers and carry no memory side-effects. They
can be replaced by equivalent sequences without additional constraints.

\subsubsection{Bitwise NOT} \label{sec:sc1-not}

The standard RISC-V pseudo-instruction \texttt{not rd, rs1} expands to \texttt{xori rd, rs1, -1}.
Neither sc0 nor sc1 include \texttt{xori}, so a different derivation is required.

\textbf{Proof.} In two's complement, negation satisfies $-x = \mathord{\sim}x + 1$ for all $x$.
Rearranging: $\mathord{\sim}x = -x - 1$. Both subtraction from zero (\texttt{sub rd, x0, rs1}) and
decrement by one (\texttt{addi rd, rd, -1}) are available in sc0 and sc1. $\square$

\begin{lstlisting}
sub  rd, x0, rs1    # rd = -rs1
addi rd, rd, -1     # rd = -rs1 - 1 = ~rs1
\end{lstlisting}

The synthesis costs 2 static instructions and requires no extra registers. It applies to both rvsc0
and rvsc1.

\subsubsection{XOR ($\mathtt{R[rd] = R[rs1] \oplus R[rs2]}$)} \label{sec:sc1-xor}

Neither sc0 nor sc1 include \texttt{xor} or \texttt{xori}.

\textbf{Proof.} Claim: $a \oplus b = (a \mid b) - (a \mathbin{\&} b)$ for all 32-bit words $a, b$,
where the subtraction is ordinary two's-complement subtraction.

Fix a bit position $i$ and write $x = a_i, y = b_i \in \{0, 1\}$. Case analysis over the four
combinations of $x, y$:

\begin{table}[htb]
\centering
\caption{Truth table establishing $(a \mid b) - (a \mathbin{\&} b) = a \oplus b$ bitwise}
\label{tab:xor-truth}
\begin{tabular}{ccccc}
\toprule
$a_i$ & $b_i$ & $(a \mid b)_i$ & $(a \mathbin{\&} b)_i$ & $(a \mid b)_i - (a \mathbin{\&} b)_i$ \\
\midrule
0 & 0 & 0 & 0 & 0 \\
0 & 1 & 1 & 0 & 1 \\
1 & 0 & 1 & 0 & 1 \\
1 & 1 & 1 & 1 & 0 \\
\bottomrule
\end{tabular}
\end{table}

The last column matches $a \oplus b$ in every row, so the identity holds per bit. Moreover
$(a \mathbin{\&} b)_i = 1$ implies $(a \mid b)_i = 1$ in every row --- the 1-bits of $a \mathbin{\&}
b$ are always a subset of the 1-bits of $a \mid b$. Consequently the per-bit subtraction never needs
to borrow from a neighboring bit position; interpreting $a \mid b$ and $a \mathbin{\&} b$ as 32-bit
binary numbers, the ordinary two's-complement subtraction \texttt{sub rd, ab\_ior, ab\_and} computes
exactly the bitwise difference, with no cross-bit borrow propagation. Hence the instruction-level
subtraction yields $a \oplus b$ exactly, for every $a, b$. $\square$

\begin{lstlisting}
and  t0, rs1, rs2   # t0 = rs1 & rs2
or   rd, rs1, rs2   # rd = rs1 | rs2
sub  rd, rd, t0     # rd = (rs1 | rs2) - (rs1 & rs2) = rs1 ^ rs2
\end{lstlisting}

The register-operand form costs 3 static instructions and 1 extra register. When the second operand
is an immediate, a \texttt{li} is prepended, raising the cost to 4 instructions and 2 extra
registers. The synthesis applies to both rvsc0 and rvsc1.

\subsubsection{Shifts: SLL, SRL, SRA} \label{sec:sc1-shifts}

Neither sc0 nor sc1 include any shift instruction. All three variants are synthesized from
\texttt{add}, \texttt{addi}, \texttt{sub}, and \texttt{beq}.

\textbf{Note on shift-count range.} The C standard (ISO C11 \S6.5.7) states that shifting a 32-bit
value by a count equal to or greater than 32 is undefined behavior; a conforming C program never
triggers this case. The synthesis sequences nonetheless mask the shift count to the range $[0, 31]$
via \texttt{and t, rs2, 31}, matching the hardware behavior of native RV32I shift instructions.

\paragraph{Logical Left Shift (SLL)} \label{sec:sc1-sll}

\texttt{sll rd, rs1, rs2} computes $\mathit{rd} = \mathit{rs1} \ll \mathit{rs2}$ (zero-fill from
the right).

\textbf{Proof by induction.} Let $\text{sll}(x, n)$ denote $x$ left-shifted by $n$ positions. Claim:
$\text{sll}(x, n) = x \cdot 2^n \pmod{2^{32}}$.

\begin{itemize}
  \item \textit{Base case} $(n = 0)$: $\text{sll}(x, 0) = x = x \cdot 2^0$. $\checkmark$
  \item \textit{Inductive step}: Assume $\text{sll}(x, n) = x \cdot 2^n$. Then
        $\text{sll}(x, n+1) = \text{sll}(x, n) + \text{sll}(x, n) = 2 \cdot x \cdot 2^n =
        x \cdot 2^{n+1}$. $\checkmark$
\end{itemize}

The synthesis implements this recursion as a count-down loop in which each iteration doubles
\texttt{rd} via \texttt{add rd, rd, rd}. $\square$

\begin{lstlisting}[language=C]
uint32_t sll(uint32_t rs1, uint32_t rs2) {
    rs2 &= 31u;           // C11 S6.5.7: mask to [0, 31]
    uint32_t rd = rs1;
    for (uint32_t i = 0; i < rs2; i++) rd += rd;
    return rd;
}
\end{lstlisting}

The resulting assembly sequence is listed in Appendix~\ref{apx:sll-asm}.

Let $b$ denote the masked shift amount. Each loop iteration performs the doubling (\texttt{add}),
decrements the counter (\texttt{addi}), and tests it (\texttt{beq}); because sc1 has no native
unconditional jump (\texttt{jal} and \texttt{auipc} are both absent), the back-edge to the top of
the loop is itself synthesized as \texttt{lui + addi + jr} (three instructions). Every iteration
therefore costs 6 instructions, except the final one, which exits through the taken \texttt{beq} and
skips the back-jump. With the mask and guard setup this gives $6b + 1$ instructions for $b \geq 1$,
a minimum of 3 when $b = 0$ and a maximum of 187 when $b = 31$, at the cost of one extra register.

\paragraph{Logical Right Shift (SRL)} \label{sec:sc1-srl}

\texttt{srl rd, rs1, rs2} computes $\mathit{rd} = \mathit{rs1} \gg \mathit{rs2}$ (zero-fill from
the left). A right shift cannot be synthesized by repeated halving, because integer division by two
would floor rather than truncate, producing incorrect results for odd values. Instead, output bit $i$
is copied from input bit $i + s$ using two single-bit masks that scan upward together.

\textbf{Proof of correctness.} The algorithm maintains \texttt{out\_mask} (scanning output bits from
0 upward) and \texttt{in\_mask} (scanning input bits from $s$ upward). At each step, if
\texttt{x \& in\_mask} $\neq$ 0, input bit $i+s$ is set and the algorithm sets
\texttt{result |= out\_mask}. Both masks then advance ($\ll\!\!= 1$). The loop terminates when
\texttt{in\_mask} overflows past bit 31 (becoming 0), meaning all input positions at or above $s$
have been processed. Output bits 0 through $31 - s$ are set from the corresponding input bits; bits
above $31 - s$ are never set (their input positions have overflowed), correctly implementing
zero-fill. $\square$

\begin{lstlisting}[language=C]
uint32_t srl(uint32_t x, uint32_t shift) {
    shift = shift & 31u;      // C11 S6.5.7: mask to [0, 31]
    uint32_t result   = 0;
    uint32_t out_mask = 1u;
    uint32_t in_mask  = 1u << shift;   // [sll]
    while (in_mask != 0) {
        if ((x & in_mask) != 0) result |= out_mask;
        out_mask <<= 1;   // [sll]
        in_mask  <<= 1;   // [sll]
    }
    return result;
}
\end{lstlisting}

The resulting assembly sequence is listed in Appendix~\ref{apx:srl-asm}. The worst-case total is
approximately 170 instructions, requiring five extra registers.

\paragraph{Arithmetic Right Shift (SRA)} \label{sec:sc1-sra}

\texttt{sra rd, rs1, rs2} computes the arithmetic right shift: identical to \texttt{srl} but vacated
upper bits are filled with the sign bit rather than zero.

\textbf{Proof.} Let $s$ be the masked shift count and $b_{31}$ the sign bit of \texttt{rs1}.
Arithmetic right shift sets output bit $i$ to input bit $\min(i+s, 31)$. For $i \leq 31-s$ this
equals the \texttt{srl} result. For $i > 31-s$, all upper bits equal $b_{31}$.

If $b_{31} = 0$, the \texttt{srl} result already has all upper bits zero; no correction is needed.

If $b_{31} = 1$, the upper $s$ bits must be 1. The mask $(-1) \ll (32 - s)$ has exactly the top $s$
bits set, since $-1$ in two's complement is all-ones, and left-shifting by $k$ clears the low $k$
bits. OR-ing this mask into the \texttt{srl} result sets the upper $s$ bits correctly. $\square$

\begin{lstlisting}[language=C]
uint32_t sra(uint32_t x, uint32_t shift) {
    shift = shift & 31u;
    uint32_t result = srl(x, shift);
    if (shift == 0 || (x >> 31) == 0) return result;
    // OR in the top 'shift' bits: left-shifting -1 leaves exactly those bits set
    uint32_t sign_mask = (uint32_t)(-1) << (32u - shift);
    return result | sign_mask;
}
\end{lstlisting}

The resulting assembly sequence is listed in Appendix~\ref{apx:sra-asm}. The synthesis reuses the
full SRL expansion and appends roughly 25 additional instructions for sign-bit extraction and
sign-mask construction, reaching a worst-case total of approximately 200 instructions at the cost of
six extra registers.

\subsection{Comparisons} \label{sec:sc1-comparisons}

\subsubsection{SLT and SLTU} \label{sec:sc1-slt}

\texttt{slt rd, rs1, rs2} sets \texttt{rd = 1} if \texttt{rs1 < rs2} (signed comparison), else
\texttt{rd = 0}. Neither sc0 nor sc1 include \texttt{slt} or \texttt{slti}.

\textbf{Proof.} Let $\mathit{diff} = \mathit{rs1} - \mathit{rs2}$ (32-bit two's complement). Without
overflow, the sign bit of $\mathit{diff}$ correctly indicates the comparison: $\mathit{diff}_{31} = 1$
iff $\mathit{rs1} < \mathit{rs2}$. Signed overflow occurs when the operands have different signs and
the result has the same sign as $\mathit{rs2}$. The expression
$\mathit{overflow} = (\mathit{rs1} \oplus \mathit{rs2}) \mathbin{\&} (\mathit{rs1} \oplus
\mathit{diff})$ has MSB 1 exactly when overflow occurred. XOR-ing $\mathit{diff}$ with
$\mathit{overflow}$ flips the sign bit iff overflow occurred, yielding the corrected result.
Verification by case analysis on the sign bits of $a = \mathit{rs1}$ and $b = \mathit{rs2}$:

\begin{itemize}
  \item $(a \geq 0, b \geq 0)$: subtraction cannot overflow; $\mathit{overflow}_{31} = 0$; result
        $= \mathit{diff}_{31}$. Correct.
  \item $(a < 0, b < 0)$: same analysis. Correct.
  \item $(a \geq 0, b < 0)$: $a \geq b$ always; result must be 0. If overflow:
        $\mathit{diff}_{31} = 1$ (wrong). $(a \oplus b)_{31} = 1$ (signs differ); $(a \oplus
        \mathit{diff})_{31} = 1$; so $\mathit{overflow}_{31} = 1$ and
        $\mathit{corrected}_{31} = 1 \oplus 1 = 0$. Correct.
  \item $(a < 0, b \geq 0)$: $a < b$ always; result must be 1. If underflow:
        $\mathit{diff}_{31} = 0$ (wrong); $\mathit{overflow}_{31} = 1$;
        $\mathit{corrected}_{31} = 0 \oplus 1 = 1$. Correct. If no underflow:
        $\mathit{diff}_{31} = 1$; $\mathit{overflow}_{31} = 0$;
        $\mathit{corrected}_{31} = 1$. Correct. $\square$
\end{itemize}

\begin{lstlisting}[language=C]
uint32_t slt(uint32_t a, uint32_t b) {
    uint32_t diff      = a - b;
    uint32_t overflow  = (a ^ b) & (a ^ diff);   // MSB = 1 iff signed overflow
    uint32_t corrected = diff ^ overflow;
    return corrected >> 31;
}
\end{lstlisting}

The resulting assembly sequence is listed in Appendix~\ref{apx:slt-asm}.

\texttt{sltu rd, rs1, rs2} performs the same comparison treating both operands as unsigned.

\textbf{Proof.} \texttt{rs1 < rs2} (unsigned) iff the subtraction \texttt{rs1 - rs2} generates a
borrow at the MSB. A borrow is \textit{generated} at bit $i$ when $\mathit{rs1}_i = 0$ and
$\mathit{rs2}_i = 1$, captured by $\mathord{\sim}\mathit{rs1} \mathbin{\&} \mathit{rs2}$. A borrow
is \textit{propagated} when $\mathit{rs1}_i = \mathit{rs2}_i$ and a borrow arrived from below,
captured by $\mathord{\sim}(\mathit{rs1} \oplus \mathit{rs2}) \mathbin{\&} \mathit{diff}$. The MSB
of their union is 1 iff \texttt{rs1 < rs2} (unsigned). $\square$

\begin{lstlisting}[language=C]
uint32_t sltu(uint32_t a, uint32_t b) {
    uint32_t diff   = a - b;
    uint32_t borrow = (~a & b) | (~(a ^ b) & diff);
    return borrow >> 31;
}
\end{lstlisting}

The resulting assembly sequence is listed in Appendix~\ref{apx:sltu-asm}.

Both expansions are expensive: \texttt{slt} requires approximately 60 instructions and three extra
registers (t0--t2), while \texttt{sltu} requires approximately 70 instructions and four extra
registers (t0--t3).

\subsubsection{BNE} \label{sec:sc1-bne}

\texttt{bne rs1, rs2, target} jumps iff \texttt{rs1} $\neq$ \texttt{rs2}.

\textbf{Equivalence.} \texttt{rs1} $\neq$ \texttt{rs2} iff \texttt{rs1 - rs2} $\neq$ 0. If the
difference is zero, \texttt{beq} falls through (skip the jump); otherwise an unconditional
\texttt{beq x0, x0, target} is taken.

\begin{lstlisting}
sub  t0, rs1, rs2
beq  t0, x0, skip    # rs1 == rs2: skip
beq  x0, x0, target  # unconditional jump
skip:
\end{lstlisting}

The expansion costs 3 instructions and one extra register.

\subsubsection{BLT, BGE, BLTU, BGEU} \label{sec:sc1-ordered-branches}

Each ordered branch is synthesized using \texttt{[slt]} or \texttt{[sltu]} followed by a \texttt{beq}
or \texttt{bne} test.

\textbf{Equivalence.} A ``branch if less-than'' is equivalent to computing the comparison into a
register and branching on that register being nonzero; ``branch if greater-than-or-equal'' branches
when the comparison is zero (i.e., the ``not less-than'' case).

\begin{table}[htb]
\centering
\caption{Ordered branch syntheses via \texttt{[slt]} and \texttt{[sltu]}}
\label{tab:ordered-branches}
\begin{tabular}{lll}
\toprule
\textbf{Branch} & \textbf{Condition} & \textbf{Synthesis} \\
\midrule
\texttt{blt rs1, rs2, L}  & \texttt{rs1 < rs2} (signed)    & \texttt{[slt t, rs1, rs2]}; \texttt{bne t, x0, L} \\
\texttt{bge rs1, rs2, L}  & \texttt{rs1 >= rs2} (signed)   & \texttt{[slt t, rs1, rs2]}; \texttt{beq t, x0, L} \\
\texttt{bltu rs1, rs2, L} & \texttt{rs1 < rs2} (unsigned)  & \texttt{[sltu t, rs1, rs2]}; \texttt{bne t, x0, L} \\
\texttt{bgeu rs1, rs2, L} & \texttt{rs1 >= rs2} (unsigned) & \texttt{[sltu t, rs1, rs2]}; \texttt{beq t, x0, L} \\
\bottomrule
\end{tabular}
\end{table}

\subsection{Immediate Variants} \label{sec:sc1-immediates}

The instructions \texttt{andi}, \texttt{ori}, \texttt{xori}, \texttt{slli}, \texttt{srli},
\texttt{srai}, \texttt{slti}, and \texttt{sltiu} encode a register together with a 12-bit signed
immediate. Neither sc0 nor sc1 support any of these forms. Each is synthesized by loading the
immediate into a temporary register via \texttt{addi} and calling the corresponding
register--register variant:

\begin{lstlisting}
# Example: andi rd, rs1, imm  ->
addi  t0, x0, imm
and   rd, rs1, t0
\end{lstlisting}

Since \texttt{addi} encodes 12-bit signed immediates (range $-2048$ to $2047$) and all of these
instruction forms share the same I-type 12-bit field, every immediate fits directly, at the cost of
one extra \texttt{addi} instruction and one extra register per immediate variant.

\subsection{Loads} \label{sec:sc1-loads}

\subsubsection{LUI (rvsc0 only)} \label{sec:sc0-lui}

\texttt{lui rd, imm20} sets $\mathit{rd} = \mathit{imm20} \ll 12$. The instruction is native in
rvsc1; synthesis is required only for rvsc0. Since rvsc0 also lacks \texttt{jalr}, function calls
are impossible; LUI synthesis is primarily needed to materialize large integer constants.

\textbf{Derivation.} A 32-bit value $v$ is split into $v = (\mathit{hi20} \ll 12) + \mathit{lo12}$,
where $\mathit{lo12}$ is the signed 12-bit remainder that makes $v - \mathit{lo12}$ an exact multiple
of 4096 (the same rounding a native \texttt{lui}+\texttt{addi} pair would use) and $\mathit{hi20}$ is
the 20-bit quantity a \texttt{lui} instruction would hold. If $\mathit{hi20}$ itself fits in
\texttt{addi} (rare but cheap when it does), it is loaded directly; otherwise, since \texttt{addi}
encodes only 12-bit signed immediates (range $-2048$ to $2047$), $\mathit{hi20}$'s low 20 bits are
split into two 10-bit halves $\mathit{hi} = \mathit{hi20}[19{:}10]$ and $\mathit{lo} =
\mathit{hi20}[9{:}0]$ (each in $[0, 1023]$, so each fits \texttt{addi}). Because these two fields
occupy disjoint bit ranges, $\mathit{hi20} = (\mathit{hi} \ll 10) + \mathit{lo}$ with no carry
between them, so the combination can use \texttt{add}/\texttt{addi} instead of \texttt{or} ---
matching the rest of this target's arithmetic-only style and needing only one register throughout:

\begin{lstlisting}
addi  rd, x0, hi     # rd = upper 10 bits of hi20  (fits in addi: 0-1023)
[sll  rd, rd, 10]    # rd = hi << 10
addi  rd, rd, lo     # rd = (hi << 10) + lo = hi20  (lo fits in addi: 0-1023)
[sll  rd, rd, 12]    # rd = hi20 << 12
addi  rd, rd, lo12   # rd = v   (only emitted if lo12 != 0)
\end{lstlisting}

Every step reuses \texttt{rd} as both source and destination, so this synthesis needs zero extra
registers and is safe to emit both before and after register allocation. Worst case (both the 10/10
split and the trailing \texttt{lo12} addition are needed) this is 25 instructions; when
$\mathit{hi20}$ itself fits \texttt{addi}, it collapses to 13--14.

An earlier version of this target instead stored the 32-bit value in a linker-placed constant pool
and loaded it with a single \texttt{lw rd, pool\_entry(x0)}, reasoning that the pool would always
link within the 12-bit signed offset of \texttt{x0} (address below 2048). That reasoning was a
mistake: nothing guarantees a program executes starting at address 0, and it did not hold even for
this project's own Spike-based behavioral tests, whose bare-metal binaries load at
\texttt{0x80000000} --- \texttt{\%lo(pool\_entry)} there computes a wrapped, incorrect address, so
the load would silently return garbage. The addi/shift synthesis above has no such dependency on
load address and was adopted instead.

\subsubsection{LB, LBU, LH, LHU} \label{sec:sc1-lb-synthesis}

Since sc1 supports only \texttt{lw} (32-bit word loads), every byte or halfword load is synthesized
in four steps: align the address to a word boundary, load the word, extract the target unit by shift,
and sign-extend or zero-extend.

\textbf{Proof of correctness.} The expression \texttt{addr \& \~{}3u} clears the two low-order bits,
yielding the address of the word that contains the byte or halfword at \texttt{addr}. The unit's
position within the word is \texttt{addr \& MASK} (MASK = 3 for bytes, 2 for halfwords), and its bit
offset is \texttt{(addr \& MASK) * 8}. Right-shifting the word by this amount moves the target unit
to bits $[N-1:0]$ where $N$ is 8 or 16. A subsequent shift pair of $32-N$ bits --- left then right
--- isolates the unit and performs either sign extension (arithmetic right shift, for \texttt{lb}/
\texttt{lh}) or zero extension (logical right shift, for \texttt{lbu}/\texttt{lhu}). This is correct
for any address alignment because the memory model is little-endian and the bit offset exactly encodes
the unit's position within the word. $\square$

\begin{lstlisting}[language=C]
/* Common algorithm for lb/lbu/lh/lhu */
uint32_t aligned  = addr & ~3u;           // word-aligned: addr & -4
uint32_t word     = lw(aligned);          // lw 0(aligned)
uint32_t unit_off = addr & MASK;          // MASK = 3 for byte, 2 for halfword
uint32_t bit_off  = unit_off * 8;         // least-significant bit position
uint32_t shifted  = word >> bit_off;      // [srl] extracts unit to bits [N:0]
// lbu/lhu: zero-extend via symmetric shift
uint32_t result   = (shifted << BITS) >> BITS;  // logical: [srl]
// lb/lh:  sign-extend via arithmetic shift
int32_t  result   = (int32_t)(shifted << BITS) >> BITS;  // arithmetic: [sra]
// BITS = 24 for byte (QImode), 16 for halfword (HImode)
\end{lstlisting}

The resulting assembly sequences are listed in Appendix~\ref{apx:lb-lh-asm}. Since sc1 shifts are
themselves synthesized via \texttt{add}/\texttt{beq} loops, each byte load expands to approximately
70--80 instructions.

\subsection{Stores} \label{sec:sc1-stores}

\subsubsection{SB and SH} \label{sec:sc1-sb}

A byte store (\texttt{sb}) or halfword store (\texttt{sh}) is implemented as a read-modify-write:
load the word containing the target location, clear the target bits, insert the new value, and write
back. The bit offset is computed at runtime since the base address is unknown at compile time.

\textbf{Proof of correctness.} Let \texttt{word} be the current word at the aligned address,
\texttt{mask} the bit-field covering the target unit (e.g., \texttt{0xFF << shift} for a byte), and
\texttt{val} the new value to write. The expression $(\mathtt{word} \mathbin{\&} \mathord{\sim}
\mathtt{mask}) \mid (\mathtt{val} \mathbin{\&} \mathtt{mask})$ replaces exactly the target bits with
\texttt{val} while preserving all others. For each bit $i$: if $\mathtt{mask}_i = 1$, then
$(\mathtt{word} \mathbin{\&} \mathord{\sim}\mathtt{mask})_i = 0$ and $(\mathtt{val} \mathbin{\&}
\mathtt{mask})_i = \mathtt{val}_i$, yielding $\mathtt{val}_i$; if $\mathtt{mask}_i = 0$, then
$(\mathtt{val} \mathbin{\&} \mathtt{mask})_i = 0$ and $(\mathtt{word} \mathbin{\&}
\mathord{\sim}\mathtt{mask})_i = \mathtt{word}_i$, yielding $\mathtt{word}_i$. $\square$

\begin{lstlisting}[language=C]
void sb(uint8_t *addr, uint32_t rs2) {
    uint32_t aligned_addr = (uint32_t)addr & ~3u;   // addr & -4
    uint32_t byte_pos     = (uint32_t)addr & 3u;    // runtime: 0-3
    uint32_t shift        = byte_pos * 8u;           // runtime: 0,8,16,24
    uint32_t old_word     = *(uint32_t *)aligned_addr;     // lw
    uint32_t byte_mask    = 0xFFu << shift;                // [sll] runtime shift
    uint32_t new_byte     = (rs2 & 0xFFu) << shift;        // [sll] runtime shift
    uint32_t new_word     = (old_word & ~byte_mask) | new_byte;
    *(uint32_t *)aligned_addr = new_word;                   // sw
}
\end{lstlisting}

The resulting assembly sequence is listed in Appendix~\ref{apx:sb-asm}. The expansion requires one
extra \texttt{lw} and one extra \texttt{sw}, totalling approximately 100 instructions and three extra
registers (t0--t2).

The \texttt{sh} synthesis follows the same read-modify-write pattern with \texttt{MASK = 2} and mask
constant \texttt{0xFFFF}. Since \texttt{0xFFFF} exceeds the 12-bit \texttt{addi} range, it is
materialized via \texttt{lui 0x10; addi -1}. The resulting assembly sequence is listed in
Appendix~\ref{apx:sh-asm}. The expansion likewise totals approximately 105 instructions and three
extra registers (t0--t2).

\subsection{Control Flow} \label{sec:sc1-jump}

\subsubsection{JAL} \label{sec:sc1-jal}

\texttt{jal ra, target} saves the return address (\texttt{PC + 4}) into \texttt{ra} and jumps to
\texttt{target}. rvsc1 does not support \texttt{auipc}-based PC-relative addressing; \texttt{jal}
must therefore be synthesized from \texttt{lui}, \texttt{addi}, and \texttt{jalr}.

\textbf{Equivalence.} The semantics of \texttt{jal ra, target} are: $\mathit{ra} \leftarrow
\mathit{PC} + 4$; $\mathit{PC} \leftarrow \mathit{target}$. Since \texttt{lui}+\texttt{addi} can
materialize any 32-bit absolute address, the two effects can be split: materialize the return address
into \texttt{ra} before the jump, then jump to \texttt{target} via \texttt{jalr}.

\begin{lstlisting}
# jal ra, target  ->
lui   ra,  %hi(back)       # ra[31:12] = upper bits of (PC+4)
addi  ra,  ra, %lo(back)   # ra = PC+4
lui   t0,  %hi(target)
addi  t0,  t0, %lo(target)
jalr  x0,  0(t0)           # PC <- target; ra already holds return address
back:
\end{lstlisting}

Each call site expands to 5 instructions and requires one extra register (\texttt{t0}).

\subsection{Cost Summary} \label{sec:sc1-synthesis-cost}

Table~\ref{tab:synthesis-cost} summarises the instruction and register cost for every synthesis
covered in this section. Counts reflect worst-case inputs (e.g.\ maximum shift amount, non-zero byte
position, misaligned halfword address).

\begin{table}[htb]
\centering
\caption{Worst-case instruction and register cost for each synthesis}
\label{tab:synthesis-cost}
\begin{tabular}{llcc}
\toprule
\textbf{Operation} & \textbf{Applies to} & \textbf{Worst-case instructions} & \textbf{Extra registers} \\
\midrule
\texttt{not}         & rvsc0, rvsc1 & 2    & 0 \\
\texttt{xor} (reg)   & rvsc0, rvsc1 & 3    & 1 \\
\texttt{xor} (imm)   & rvsc0, rvsc1 & 4    & 2 \\
\texttt{ori} (imm)   & rvsc1        & 2    & 1 \\
\texttt{andi} (imm)  & rvsc1        & 2    & 1 \\
\texttt{sll}         & rvsc0, rvsc1 & 127  & 1 \\
\texttt{srl}         & rvsc0, rvsc1 & $\sim$170 & 5 \\
\texttt{sra}         & rvsc0, rvsc1 & $\sim$200 & 6 \\
\texttt{slt}         & rvsc0, rvsc1 & $\sim$60  & 3 \\
\texttt{sltu}        & rvsc0, rvsc1 & $\sim$70  & 4 \\
\texttt{bne}         & rvsc0, rvsc1 & 3    & 1 \\
\texttt{blt}/\texttt{bltu} & rvsc0, rvsc1 & $\sim$75 & 4 \\
\texttt{bge}/\texttt{bgeu} & rvsc0, rvsc1 & $\sim$75 & 4 \\
\texttt{lui} (addi+shift) & rvsc0   & 25   & 0 \\
\texttt{lb}/\texttt{lbu} & rvsc0, rvsc1 & $\sim$80 & 2 \\
\texttt{lh}/\texttt{lhu} & rvsc0, rvsc1 & $\sim$80 & 2 \\
\texttt{sb}          & rvsc0, rvsc1 & $\sim$100 & 3 \\
\texttt{sh}          & rvsc0, rvsc1 & $\sim$105 & 3 \\
\texttt{jal}         & rvsc1        & 5    & 1 \\
\bottomrule
\end{tabular}
\end{table}

\section{GCC Implementation} \label{sec:sc1-gcc-impl}

\subsection{Target Flags} \label{sec:sc1-flags}

Each synthesized instruction corresponds to a Boolean flag declared in
\texttt{gcc/config/riscv/riscv.opt}. When a flag is zero, the corresponding \texttt{define\_insn}
condition evaluates to false (suppressing native instruction emission) while the
\texttt{define\_expand} synthesis path fires and calls \texttt{DONE}.
Table~\ref{tab:target-flags} lists all flags added for this project.

\begin{table}[htb]
\centering
\caption{Target option flags added for this project}
\label{tab:target-flags}
\begin{tabular}{lll}
\toprule
\textbf{Flag} & \textbf{Macro} & \textbf{Disabled for} \\
\midrule
\texttt{-mfence}  & \texttt{TARGET\_FENCE}  & rvsc0, rvsc1, rvsc2 \\
\texttt{-mauipc}  & \texttt{TARGET\_AUIPC}  & rvsc0, rvsc1 \\
\texttt{-mlui}    & \texttt{TARGET\_LUI}    & rvsc0 \\
\texttt{-mshift}  & \texttt{TARGET\_SHIFT}  & rvsc0, rvsc1 \\
\texttt{-mxor}    & \texttt{TARGET\_XOR}    & rvsc0, rvsc1 \\
\texttt{-mori}    & \texttt{TARGET\_ORI}    & rvsc0, rvsc1 \\
\texttt{-mandi}   & \texttt{TARGET\_ANDI}   & rvsc0, rvsc1 \\
\texttt{-mbne}    & \texttt{TARGET\_BNE}    & rvsc0, rvsc1 \\
\texttt{-mslt}    & \texttt{TARGET\_SLT}    & rvsc0, rvsc1 \\
\texttt{-mslti}   & \texttt{TARGET\_SLTI}   & rvsc0, rvsc1 \\
\texttt{-mblt}    & \texttt{TARGET\_BLT}    & rvsc0, rvsc1 \\
\texttt{-mbge}    & \texttt{TARGET\_BGE}    & rvsc0, rvsc1 \\
\texttt{-mbltu}   & \texttt{TARGET\_BLTU}   & rvsc0, rvsc1 \\
\texttt{-mbgeu}   & \texttt{TARGET\_BGEU}   & rvsc0, rvsc1 \\
\texttt{-mbyte}   & \texttt{TARGET\_BYTE}   & rvsc0, rvsc1 \\
\texttt{-mhalf}   & \texttt{TARGET\_HALF}   & rvsc0, rvsc1 \\
\bottomrule
\end{tabular}
\end{table}

All flags have \texttt{Init(1)} (enabled by default). The per-target header (\texttt{rvscN.h})
disables the appropriate subset via \texttt{CC1\_SPEC}.

\subsection{Machine Description Patterns} \label{sec:sc1-md}

The GCC machine description (\texttt{riscv.md}) uses two complementary constructs per synthesized
operation: a \texttt{define\_insn} guarded by \texttt{TARGET\_XYZ} that emits the native instruction
when the flag is enabled, and a \texttt{define\_expand} guarded by \texttt{!TARGET\_XYZ} that emits
the synthesis sequence and calls \texttt{DONE}, preventing GCC from falling through to the native
insn. The XOR synthesis illustrates the pattern:

\begin{lstlisting}[language=Lisp]
;; Native XOR (and OR) -- emitted only when TARGET_XOR (XOR) or always (IOR)
(define_insn "*<optab><mode>3"
  [(set (match_operand:X                0 "register_operand" "=r,r")
        (any_or:X (match_operand:X 1 "register_operand" "%r,r")
                       (match_operand:X 2 "arith_operand"    " r,I")))]
  "TARGET_XOR || (<CODE>) == IOR"
  "<insn>%i2\t%0,%1,%2")

;; XOR/OR synthesis -- the any_or expand shared by both codes
(define_expand "<optab><mode>3"
  [(set (match_operand:X 0 "register_operand")
        (any_or:X (match_operand:X 1 "register_operand" "")
                   (match_operand:X 2 "reg_or_const_int_operand" "")))]
  ""
{
  /* sc1 synthesis: a ^ b = (a | b) - (a & b). */
  if ((<CODE>) == XOR && !TARGET_XOR)
    {
      rtx op1    = operands[1];
      rtx op2    = REG_P (operands[2]) ? operands[2]
                                       : force_reg (SImode, operands[2]);
      rtx ab_and = gen_reg_rtx (SImode);
      rtx ab_ior = gen_reg_rtx (SImode);
      emit_insn (gen_andsi3 (ab_and, op1, op2));
      emit_insn (gen_iorsi3 (ab_ior, op1, op2));
      emit_insn (gen_subsi3 (operands[0], ab_ior, ab_and));
      DONE;
    }
  /* sc1 synthesis: ori rd, rs, imm -> li t, imm; or rd, rs, t */
  if ((<CODE>) == IOR && !TARGET_ORI && CONST_INT_P (operands[2]))
    operands[2] = force_reg (<MODE>mode, operands[2]);
  if (CONST_INT_P (operands[2]) && synthesize_ior_xor (<OPTAB>, operands))
    DONE;
})
\end{lstlisting}

The pattern is shared between \texttt{XOR} and \texttt{IOR} via the \texttt{any\_or} code iterator.
The \texttt{DONE} call inside the \texttt{XOR} branch signals that the expansion body has fully
handled the operation; GCC does not attempt to match the \texttt{define\_insn} afterwards.

\subsection{Worked Example: XOR in C to Assembly} \label{sec:sc1-example}

The following traces how \texttt{unsigned f(unsigned a, unsigned b) \{ return a \^{} b; \}} is
compiled by \texttt{rvsc1-unknown-elf-gcc -S -O1}:

\begin{enumerate}
  \item \textbf{Frontend}: parses \texttt{a \^{} b} to an AST \texttt{XOR\_EXPR} node.
  \item \textbf{GIMPLE}: \texttt{\_1 = a \^{} b; return \_1;} --- language-independent SSA form.
  \item \textbf{RTL lowering}: GCC attempts to emit \texttt{xorsi3}. \texttt{TARGET\_XOR} is 0
        (rvsc1 disables XOR), so the \texttt{XOR} branch of the synthesis body fires
        (Section~\ref{sec:sc1-md}).
  \item \textbf{Expansion}: the body allocates two pseudo-registers (\texttt{ab\_and},
        \texttt{ab\_ior}), emits \texttt{andsi3} into \texttt{ab\_and}, \texttt{iorsi3} into
        \texttt{ab\_ior}, then \texttt{subsi3} computing \texttt{ab\_ior - ab\_and} into the
        destination, and calls \texttt{DONE}.
  \item \textbf{Register allocation}: GCC maps pseudo-registers to physical registers (\texttt{a0},
        \texttt{a1}, \texttt{a5}) such that no \texttt{xor} instruction appears; \texttt{ab\_ior} is
        allocated directly into the destination register \texttt{a0}, so only one extra register
        (\texttt{a5}, holding \texttt{ab\_and}) is visible in the output.
  \item \textbf{Assembly output}:
\end{enumerate}

\begin{lstlisting}
and  a5, a0, a1
or   a0, a0, a1
sub  a0, a0, a5
\end{lstlisting}

Inspecting the result with \texttt{grep xor} returns empty; only \texttt{and}, \texttt{or}, and
\texttt{sub} appear --- all native sc1 instructions.

\section{Known Limitations}

\subsection{rvsc0: No Non-Inlined Function Calls}

The rvsc0 processor supports neither \texttt{jalr} nor \texttt{jal}. The compiler therefore has no
instruction with which to perform an indirect jump while saving the return address, making non-inlined
function calls impossible. However, the GCC inliner operates before instruction selection, so C
programs with multiple functions can still be compiled for rvsc0 provided that every call site is
inlined, either by marking functions \texttt{\_\_attribute\_\_((always\_inline))} or by relying on
GCC's automatic inlining at \texttt{-O1} and above. Any call that remains non-inlined after
optimization will either fail at link time or produce incorrect control flow at runtime.

\subsection{rvsc1: Function Body Size Limit}

The only native conditional branch in rvsc1 is \texttt{beq}, which has a $\pm$4~KB PC-relative
offset range (12-bit signed immediate). For branch targets that exceed this range, the GCC backend
emits a long form:

\begin{lstlisting}
beq  rs1, rs2, skip        # short branch: skip the jump if condition holds
lui  t1, %hi(target)
addi t1, t1, %lo(target)
jr   t1                    # absolute jump to target
skip:
\end{lstlisting}

This long form is itself valid rvsc1 code, so correctness is preserved. However, the long form for
the inverted branch (\texttt{bne}) chains through the \texttt{bne} synthesis (which uses
\texttt{beq}), so individual function bodies should remain under approximately 4~KB of machine code,
roughly 1\,000 synthesized instructions, to avoid triggering branch relaxation in unexpected cases.
Typical educational programs are well within this limit.
